\documentclass[11pt]{article}
\usepackage{amsmath,amssymb,amsthm}

\newtheorem{theorem}{Theorem}
\newtheorem{corollary}[theorem]{Corollary}
\newtheorem{lemma}[theorem]{Lemma}

\newcommand{\Up}{U_+}
\newcommand{\mass}{\operatorname{B}}

\title{The complete excess-two inverse decoder layer}
\author{}
\date{27 August 2026}

\begin{document}
\maketitle

For a positive odd integer $x$, let
\[
 \Up(x):=\frac{3x+1}{2^{v_2(3x+1)}}.
 \tag{1}
\]
An exponent word $E=(e_1,\ldots,e_h)$ is read from its positive odd
start toward its terminal state.  Its mass and excess are
\[
 \mass(E):=e_1+\cdots+e_h,
 \qquad t(E):=\mass(E)-h.
 \tag{2}
\]
This note classifies every exact inverse word with $t(E)=2$ at a fixed
positive odd terminal state.

The elementary reverse rule is useful.  If $x$ is the current odd
terminal state, an exponent-$a$ predecessor is
\[
 p=\frac{2^a x-1}{3}.
 \tag{3}
\]
It exists precisely when $3\nmid x$ and
\[
 a\equiv
 \begin{cases}
 1\pmod2,&x\equiv2\pmod3,\\
 0\pmod2,&x\equiv1\pmod3.
 \end{cases}
 \tag{4}
\]
In shifted coordinates $X=x+1$, reversing exponents one, two, and three
uses respectively
\[
 X\longmapsto\frac{2X}{3},
 \qquad
 X\longmapsto\frac{4X-2}{3},
 \qquad
 X\longmapsto\frac{8X-6}{3}.
 \tag{5}
\]

\begin{theorem}[complete excess-two classification]
Let $r$ be a positive odd integer and put
\[
 b:=v_3(r+1),
 \qquad s:=\frac{r+1}{3^b}.
 \tag{6}
\]
There are exactly two families of exact predecessor words $E$ with
$t(E)=2$.

\medskip
\noindent\textbf{Family I: one exponent three.}
For every $k$ with $0\leq k<b$, define
\[
 M_k:=2^{k+2}3^{b-k-1}s-1,
 \qquad c_k:=v_3(M_k).
 \tag{7}
\]
For every $0\leq\ell\leq c_k$, there is one exact word
\[
 \boxed{
 E=1^\ell3\,1^k,
 \qquad |E|=k+1+\ell,
 \qquad z+1=2\left(\frac23\right)^\ell M_k.}
 \tag{8}
\]
These are all the excess-two words containing an exponent three.

\medskip
\noindent\textbf{Family II: two exponents two.}
Define
\[
 N_1:=2^{b+1}s-1,
 \qquad b_1:=v_3(N_1).
 \tag{9}
\]
This family is empty if $b_1=0$.  If $b_1\geq1$, put
\[
 s_1:=\frac{N_1}{3^{b_1}},
 \qquad N_2:=2^{b_1+1}s_1-1,
 \qquad b_2:=v_3(N_2).
 \tag{10}
\]
For every $1\leq j\leq b_2$, there is one exact word
\[
 \boxed{
 E=1^{j-1}2\,1^{b_1-1}2\,1^b,
 \qquad |E|=b+b_1+j,
 \qquad z+1=\left(\frac23\right)^jN_2.}
 \tag{11}
\]
These are all the excess-two words containing two exponents two.
There are no other exact words of excess two.
\end{theorem}

\begin{proof}
Because all exponents are positive and their excesses sum to two, the
word has exactly one of the two forms indicated: one exponent three and
all other exponents one, or two exponents two and all other exponents
one.

Read the word backward from $r$, and let $k$ be the number of trailing
exponents one.  Rule (5) shows that precisely $b$ consecutive
exponent-one inverses are available at the outset.

First suppose that the next non-one exponent is three.  Its parity in
(4) is odd, so it must occur before the exponent-one run is exhausted;
thus $0\leq k<b$.  After reversing the $k$ trailing ones, the shifted
state is
\[
 X=2^k3^{b-k}s.
 \tag{12}
\]
Reversing exponent three gives
\[
 \frac{8X-6}{3}
 =2\bigl(2^{k+2}3^{b-k-1}s-1\bigr)
 =2M_k.
 \tag{13}
\]
All remaining exponents must equal one.  Exactly $c_k=v_3(M_k)$ of
them are available, and $\ell$ such steps give shifted start
\[
 \left(\frac23\right)^\ell2M_k.
\]
This proves (7)--(8), both existence and exhaustion of Family I.

Now suppose that the word contains two exponents two.  Even parity in
(4) forces the first such exponent, read backward, to occur only after
all $b$ initial exponent-one inverses.  At that point the shifted state
is $2^bs$, and reversing exponent two would give
\[
 \frac{4\,2^bs-2}{3}=\frac{2N_1}{3}.
 \tag{14}
\]
It is possible exactly when $b_1\geq1$.  The shifted state in (14) has
$3$-adic valuation $b_1-1$.  Since the only remaining unit of excess
must form another exponent two, parity forces the intervening run to
have exactly $b_1-1$ exponents one.  After that run the shifted state is
\[
 2^{b_1}s_1.
 \tag{15}
\]
The second exponent-two inverse gives $2N_2/3$ and is possible exactly
when $b_2\geq1$.  Thereafter all exponents equal one; reversing $j-1$
of them is possible precisely for $1\leq j\leq b_2$, and gives shifted
start
\[
 \left(\frac23\right)^{j-1}\frac{2N_2}{3}
 =\left(\frac23\right)^jN_2.
 \tag{16}
\]
The word and length in (11) follow.  The parity alternatives above are
forced at every first non-one position, so the two families exhaust all
excess-two words.
\end{proof}

\begin{corollary}[finite excess-two horizon]
With the notation of Theorem~1, set
\[
 H_I:=\max_{0\leq k<b}(k+1+c_k),
 \tag{17}
\]
where the maximum of the empty set is zero, and set
\[
 H_{II}:=
 \begin{cases}
 b+b_1+b_2,&b_1\geq1\text{ and }b_2\geq1,\\
 0,&b_1b_2=0.
 \end{cases}
 \tag{18}
\]
Then no exact $h$-step predecessor word of $r$ has excess two when
\[
 \boxed{h>H_2(r):=\max(H_I,H_{II}).}
 \tag{19}
\]
Combining this with the excess-zero and excess-one decoder theorem, put
\[
 H_{\leq2}(r):=\max\bigl(b,\ b+b_1,\ H_2(r)\bigr).
 \tag{20}
\]
Every exact $h$-step predecessor word with $h>H_{\leq2}(r)$ has excess
at least three.
\end{corollary}

\begin{proof}
The largest length in Family I is (17), and the largest possible length
in Family II is (18).  This proves (19).  Excess zero ends at depth $b$,
and excess one ends at depth $b+b_1$.  Taking the maximum proves (20).
\end{proof}

\begin{corollary}[fixed-seed inequalities]
Let $n$ be the globally least positive odd integer with an unbounded
accelerated orbit, and suppose the orbit of $r$ is unbounded.  Then every
Family I branch satisfies
\[
 \boxed{
 M_k\geq\frac{n+1}{2}\left(\frac32\right)^\ell
 \qquad(0\leq\ell\leq c_k),}
 \tag{21}
\]
and every Family II branch satisfies
\[
 \boxed{
 N_2\geq(n+1)\left(\frac32\right)^j
 \qquad(1\leq j\leq b_2).}
 \tag{22}
\]
In particular, at the least same-depth suffix $q_*$ of a record-atom
block, all variable-depth prefixes of excess at most two are now given
by finite explicit valuation ladders.  Beyond $H_{\leq2}(q_*)$, the
minimum suffix channel has prefix excess at least three.
\end{corollary}

\begin{proof}
Every positive predecessor of $r$ also has an unbounded forward orbit,
so its start $z$ is at least $n$.  Substitution of the two exact shifted
starts in (8) and (11) gives (21) and (22).  The final assertion is
Corollary~2 applied to $r=q_*$.
\end{proof}

\end{document}
